\documentclass[11pt,a4paper]{article}

\usepackage[margin=2.2cm]{geometry}
\usepackage{booktabs}
\usepackage{array}
\usepackage{tabularx}
\usepackage{amsmath}
\usepackage{amssymb}
\usepackage[hidelinks]{hyperref}
\usepackage{enumitem}

\newcolumntype{L}{>{\raggedright\arraybackslash}p{0.16\textwidth}}
\newcolumntype{C}{>{\raggedright\arraybackslash}p{0.30\textwidth}}
\newcolumntype{R}{>{\raggedright\arraybackslash}p{0.44\textwidth}}

\newcommand{\code}[1]{\texttt{#1}}

\title{Advection--Diffusion (Advdif) Case:\\
A Comparison of \code{phys-vae} and the \code{interpretable-diffusion} PhysVAE Reimplementation\\
\large With a ranked diagnosis of the coefficient-estimation accuracy gap}
\author{Generated by Mistral Vibe}
\date{\today}

\begin{document}
\maketitle

\section*{Scope and sources}

This report compares the advection--diffusion (advdif) experiment of the original
\code{phys-vae} repository (Takeishi \& Kalousis, \emph{Physics-Integrated Variational
Autoencoders}, NeurIPS 2021; files \code{physvae/advdif/*}, \code{data/advdif/*},
\code{advdif\_train.sh}, \code{advdif\_extpol.ipynb}) against the PhysVAE
reimplementation in the private \code{interpretable-diffusion} repository
(files \code{src/models/phys\_vae.py}, \code{src/models/train\_phys\_vae.py},
\code{src/datasets/pivae.py}, \code{src/models/physics\_models.py},
\code{src/train.py}, \code{src/eval.py}, and the configuration
\code{configs/diffusion\_vae/26MAR30.yaml}).

The reimplementation currently does not match the original's diffusion-coefficient
estimation accuracy. Section~\ref{sec:diagnosis} ranks the causes found by static
analysis of both codebases. No training runs were executed in the environment where
this report was produced; the ranking is code-based, and each item states how it can
be confirmed empirically.

\section{Setup at a glance}
\label{sec:setup}

\begin{table}[h]
\centering
\small
\begin{tabularx}{\textwidth}{LXX}
\toprule
 & \textbf{Original (\code{phys-vae})} & \textbf{Reimplementation} \\
\midrule
Framework & PyTorch + torchdiffeq & JAX / Flax nnx + diffrax \\
Task & $y_t = a\,y_{xx} - b\,y_x$; known physics = diffusion only ($a$ = \code{dcoeff}); advection ($b$) handled by an in-ODE neural network & Same concept: known physics = $a\,\partial^2 y/\partial s^2$; advection $b$ handled by a \code{pre\_func} network \\
Latents (\code{physnn} recipe) & $z_\mathrm{phy} = a$ (1-D), $z_{\mathrm{aux1}} \in \mathbb{R}^4$ (in-ODE), no out-ODE auxiliary & Phys latent $=a$ (1-D), one pre-func with latent dim 4, no post-funcs \\
Data shape & $(N, 20, 50)$ = (samples, space, time) & $(N, 12, 50)$, permuted/flattened by \code{PiVAEDataset} \\
\bottomrule
\end{tabularx}
\end{table}

\section{Data generation}
\label{sec:data}

\begin{table}[h]
\centering
\small
\begin{tabularx}{\textwidth}{LXX}
\toprule
\textbf{Aspect} & \textbf{Original} & \textbf{Reimplementation} \\
\midrule
Spatial grid & 20 points, $dx = 0.1$ (\code{data/advdif/makedata.sh}) & 12 points, $X = 2.0$, $ds = 2/11 \approx 0.182$ (\code{pivae.py:198--203}) \\
Time grid & 50 frames, $dt = 0.02$ (span 0.98); \emph{$dt$ is stored in \code{args\_*.json} and read by the trainer} (\code{physvae/advdif/train.py:69--70}) & 50 frames, $T = 1.0$, so $dt = 1/49 \approx 0.0204$ (span 1.0) (\code{pivae.py:198}) \\
PDE solver & \code{scipy.solve\_ivp}, DOP853, essentially exact (\code{generate.py:12--18}) & \code{diffrax} Euler at $dt = 1/49$ (\code{pivae.py:217--221}) \\
Discretization & 3-point central differences via \code{convolve1d}: $y_{xx} = [1,-2,1]/dx^2$ (\code{generate.py:13--15}) & \code{jnp.gradient} applied twice (\code{pivae.py:131--136}); interior Laplacian is the wide 5-point stencil $(y_{i+2}-2y_i+y_{i-2})/(4ds^2)$ \\
Boundary & Zero ghost cells; boundary values evolve & $\partial y/\partial t$ padded with 0 at the edges, so boundary values are frozen (\code{pivae.py:138}, \code{physics\_models.py:66--70}) \\
Coefficients & $a, b \sim U[10^{-2}, 10^{-1}]$; initial amplitude $c \sim U[0.5, 1.5]$; $c\sin(\pi s/L)$ & $a, b \sim U[0.01, 0.1]$; $c \sim U[0.5,1.5]$; $c\sin(\pi s/X)$ \\
Observation noise & $\sigma = 10^{-3}$ & $\sigma = 10^{-3}$ (\code{pivae.py:225}) \\
Splits & train 1000 / valid 500 / test 1000, separate seeds & 3500 generated: test 500 / val 1000 / pool 2000 with 1000 used (\code{pivae.py:373--380}) \\
Saved targets & \code{init\_mag}, \code{dcoeff}, \code{ccoeff} all saved (\code{generate.py:96--98}) & Only $a$ is returned as the coefficient target (\code{pivae.py:225--227}) \\
\bottomrule
\end{tabularx}
\end{table}

\noindent Key observations:
\begin{itemize}[nosep]
  \item The reimplementation's grid is coarser (12 vs 20 points) and its effective
        Laplacian is a wide, smoothed stencil, which weakens the curvature signal from
        which $a$ must be inferred. Both are internally consistent between the
        reimplementation's data generator and its physics model, so they are deviations
        from the original rather than self-inconsistencies.
  \item The original \emph{cannot} suffer a model/data time-grid mismatch because
        $dt$ flows from the data-generation JSON into the trainer. The reimplementation
        decouples them (config \code{timestep} vs generator $T$), which enables the
        most damaging error in Section~\ref{sec:diagnosis} (D1).
\end{itemize}

\section{Architectures}
\label{sec:arch}

\begin{table}[h]
\centering
\small
\begin{tabularx}{\textwidth}{LXX}
\toprule
\textbf{Component} & \textbf{Original} & \textbf{Reimplementation} \\
\midrule
Aux encoder & Dedicated feature extractor (MLP $[256,256] \to 256$) per latent, then small heads (\code{model.py:62--80}) & Monolithic MLP $[256,256,256,64,32] \to 2\cdot\text{latent}$, split into $(\mu, \log\sigma)$ (\code{phys\_vae.py:104--110}) \\
Physics-encoder path & Unmixer (linear map + residual MLP from auxiliary means) $\to$ separate feature extractor $\to$ \code{dcoeff} head $=$ MLP $+$ \textbf{Softplus} (\code{model.py:94--101}) & Unmixer (Linear(12,12) + residual MLP, optional skip) $\to$ phys MLP $\to$ split $(\mu,\log\sigma)$; the $\mu$ head is \textbf{linear and can be negative} (\code{phys\_vae.py:128--166}) \\
Positivity / clipping & Softplus head; after sampling, clamp to feasible range $[5\cdot10^{-3},\,2\cdot10^{-1}]$, wider than the data range (\code{model.py:18, 341--342}) & No softplus; sample clipped to the data range $[0.01, 0.1]$ (\code{phys\_vae.py:668--676}); zero gradient outside the range \\
Decoder mixing & Learnable linear map $y \mapsto x$ (\code{func\_aux2\_map}, \code{model.py:44--45}), active even with the out-ODE auxiliary disabled & Identity: $\hat{x}$ is the ODE solution directly (\code{phys\_vae.py:478--486}) \\
In-ODE auxiliary & MLP($[z_{\mathrm{aux1}}, y, \sin t, \cos t] \to [64,64] \to y$), input dim $4+20+2$ & Same form: \code{pre\_func}($[z, y, \sin t, \cos t] \to [64,64] \to 12$) (\code{phys\_vae.py:238--247}) --- faithful \\
Physics model & $a \cdot \mathrm{conv1d}([1,-2,1])/dx^2$ (\code{model.py:143--158}) & $a \cdot \nabla^2_s / ds^2$ via double \code{jnp.gradient}, zero-padded edges (\code{physics\_models.py:23--70}) \\
Integration & \code{odeint} over \code{t\_intg}, $dt_\mathrm{intg} = dt/\mathrm{intg\_lev}$, \textbf{intg\_lev $=$ 2}; $dt$ taken from the data file & diffrax Euler, \code{integral\_level} $=1$, $dt_0 =$ config \code{timestep} (\code{phys\_vae.py:525--546}); also integrates to \code{t1 = integral\_len} (a step \emph{count}, not a time), so each decode integrates to $t=50$ while only saving up to $t = 2.45$ (\code{phys\_vae.py:345--350}) --- roughly $20\times$ wasted compute per decode \\
Combination enumeration & Hand-written $x_P, x_{PA}, x_{PB}, x_{PAB}$ decodes (\code{model.py:213--289}) & Generalized bitmask enumeration of pre/post-func combinations (\code{gen\_selections}, \code{full\_forward}) --- a genuine generalization, but see bug D5 \\
\bottomrule
\end{tabularx}
\end{table}

\section{Training losses and objectives}
\label{sec:losses}

Both implement the same paper: an ELBO plus three regularizers. The correspondence and
the weights used in the original's advdif \code{physnn} recipe (\code{advdif\_train.sh})
versus the reimplementation's config (\code{configs/diffusion\_vae/26MAR30.yaml}):

\begin{table}[h]
\centering
\small
\begin{tabularx}{\textwidth}{LXX}
\toprule
\textbf{Term} & \textbf{Original} & \textbf{Reimplementation} \\
\midrule
Reconstruction & $\sum$ over all $20\times50$ dims of $(\hat{x}-x)^2$, mean over batch (\code{train.py:71}) & $\sum$ over \textbf{time only}, then mean over batch \emph{and space} (\code{train\_phys\_vae.py:86}) --- about $12\times$ weaker in scale \\
KL balance & $kld \cdot e^{x\_\mathrm{lnvar}}$ with $x\_\mathrm{lnvar} = -13.8$, i.e.\ weight $\approx 10^{-6}$ (\code{train.py:109,144}) & $kld \cdot e^{x\_\mathrm{log\_\sigma}}$ with $x\_\mathrm{log\_\sigma} = 0.01$, i.e.\ weight $\approx 1.01$ (\code{train\_phys\_vae.py:222}, config line 30) --- \textbf{$\sim10^6\times$ stronger KL} \\
Priors & \code{dcoeff} prior range $[5\cdot10^{-3},\,2\cdot10^{-1}]$, deliberately wider than the data range $[10^{-2},10^{-1}]$; prior mean $0.1025$ lies \emph{outside} the data range; prior std $=0.866\times$ range (\code{model.py:181--186}) & Phys prior range $[0.01, 0.1]$ = the data range exactly; prior mean $0.055$ (\code{phys\_vae.py:587--600}, config line 34). Same $0.866$ formula, but a far more informative prior \\
Unmixing reg.\ ($R_{\mathrm{DA},1}$) & weight $\gamma = 10^{-2}$: $\|z_\mathrm{unmixed} - x_P.\mathrm{detach}()\|^2$ (\code{train.py:129--131}) & \code{physics\_only\_loss}, weight $\beta = 0.01$, with \code{stop\_gradient} (\code{train\_phys\_vae.py:94--114}) --- faithful \\
Data-augmentation reg.\ ($R_{\mathrm{DA},2}$) --- \textbf{the only direct supervision of $a$} & weight $\beta = 10^{6}$: decode pure-physics data with random \code{dcoeff}, feed through feature + \code{dcoeff} head, MSE (\code{train.py:133--140}) & weight $\gamma = 1$ (\code{augmentation\_loss}, \code{train\_phys\_vae.py:116--163}) --- \textbf{$10^6\times$ weaker} \\
Least-action / discrepancy ($R_\mathrm{ppc}$) & weight $\alpha = 10^{-1}$; $\tfrac14(\|PA-P\|^2 + \|PB-P\|^2 + \|PAB-PA\|^2 + \|PAB-PB\|^2)$ (\code{train.py:142--149}); $\alpha$ also added to the KL weight (total 1.1) & \code{discrepancy\_loss} over Hamming-distance-1 combination pairs, weight $\alpha = 0.1$ (\code{train\_phys\_vae.py:166--196}) --- right idea, broken implementation (bug D5) \\
ELBO target & $x_{PAB}$ (physics + all auxiliary) & \code{x\_stat[0][-1]} (last combination = all auxiliary applied) --- faithful \\
\bottomrule
\end{tabularx}
\end{table}

\section{Training protocol}
\label{sec:protocol}

\begin{table}[h]
\centering
\small
\begin{tabularx}{\textwidth}{LXX}
\toprule
\textbf{Aspect} & \textbf{Original} & \textbf{Reimplementation} \\
\midrule
Epochs & 20{,}000 ($\to$ 100k steps at 5 batches/epoch) & \textbf{100} ($\to$ 500 steps) (config line 50) \\
Optimizer & Adam, lr $10^{-3}$, $\varepsilon = 10^{-3}$, \textbf{weight decay $10^{-6}$}, \textbf{value grad-clip 5.0} & Adam, lr $10^{-3}$, $\varepsilon = 10^{-3}$; no weight decay, no clipping \\
Validation & Pure no-grad evaluation; \textbf{best-validation-reconstruction checkpoint saved} (\code{train.py:245--250}) & Validation loop \textbf{calls \code{train\_step}, i.e.\ runs optimizer updates on the validation data} (\code{src/train.py:133}); the final model is saved, no best-checkpoint selection (\code{main.py:39}) \\
Batch size / train size & 200 / 1000 & 200 / 1000 --- match \\
\bottomrule
\end{tabularx}
\end{table}

\section{Evaluation}
\label{sec:eval}

\begin{table}[h]
\centering
\small
\begin{tabularx}{\textwidth}{LXX}
\toprule
\textbf{Aspect} & \textbf{Original} & \textbf{Reimplementation} \\
\midrule
Coefficient metric & Paper Table 1 reports the inference error of \code{dcoeff} on the test split, using the posterior statistic & \code{coeff\_mae} $=$ mean absolute error of the predicted $a$ vs true $a$ (\code{eval.py:82}) --- absolute MAE, not the paper's error definition, so the numbers are not directly comparable \\
Estimate used & Posterior mean of $q(a \mid x)$ & \textbf{A reparameterized sample} ($\mu + \varepsilon e^{\log\sigma}$, then clipped) via \code{PhysVAE.\_\_call\_\_} (\code{phys\_vae.py:566--577}, \code{eval.py:74}) --- sampling noise inflates the MAE, severely if $\log\sigma$ is pushed up (bug D4) \\
Other metrics & Extrapolation visualization to 1000 steps (\code{advdif\_extpol.ipynb}) & Forecast MAE over all frames \\
Reference numbers & Per the paper's Table 1 (as summarized in later surveys of it): NN+phys+reg inference error for $a \approx 0.0095$; removing the data-augmentation regularizer alone degrades it to $\approx 0.199$ ($\sim20\times$) & --- \\
\bottomrule
\end{tabularx}
\end{table}

\section{Ranked diagnosis of the coefficient-accuracy gap}
\label{sec:diagnosis}

\subsection*{D1 --- Time-grid mismatch between decoder and data (most likely the dominant error)}

The generator uses $T = 1.0$ with 50 frames, so $dt = 1/49 \approx 0.0204$
(\code{pivae.py:198}). The model config uses \code{timestep: 0.05} and
\code{integral\_level: 1}, so the decoder's 50 frames sit at
$t = 0, 0.05, \dots, 2.45$ (\code{phys\_vae.py:525--530}). The decoder's time axis is
stretched $2.45\times$ relative to the data's. For the fit to succeed, the effective
per-frame diffusion rate must satisfy $a \cdot 0.05 \approx a_\mathrm{true} \cdot 0.0204$,
so either the physics term carries the fit and $a$ is systematically underestimated by
$\approx 2.45\times$ (a $\sim59\%$ relative bias, which alone gives an MAE around
$0.03$ for $a \in [0.01, 0.1]$), or the auxiliary network absorbs the difference and
$a$ becomes unidentifiable. The original avoids this by construction: it reads $dt$
from the data's \code{args.json} (\code{physvae/advdif/train.py:69--70}).

\textbf{Fix:} set \code{timestep} to $1/49$ (and preferably \code{integral\_level: 2}
like the original), or regenerate the data with $dt = 0.05$ and $T = 2.45$; better,
derive \code{timestep} from the generator the way the original does.

\subsection*{D2 --- KL term effectively a million times stronger}

The original balances the KL by the observation variance $e^{-13.8} \approx 10^{-6}$,
so the ELBO barely regularizes and the encoder can learn sharp, informative
posteriors. The reimplementation balances by $e^{0.01} \approx 1.01$
(\code{train\_phys\_vae.py:222}). With KL at full strength against a weak
reconstruction term (D8), the posterior over $a$ is pulled toward the prior ---
mean $0.055$ --- which yields an MAE of $\mathbb{E}|a - 0.055| \approx 0.025$: exactly
the collapsed regime.

\textbf{Fix:} \code{x\_log\_sigma} must encode the data noise scale: $e^{x\_\mathrm{log\_\sigma}}$
should be $\approx \sigma^2 = 10^{-6}$, i.e.\ \code{x\_log\_sigma} $\approx -13.8$,
restoring the original's semantics.

\subsection*{D3 --- Data-augmentation regularizer weight 1 vs $10^6$}

This regularizer is the only term that directly supervises the $a$-encoder on
ground-truth coefficients. The paper's own ablation shows coefficient inference
collapses without it. In the reimplementation it is $10^6\times$ weaker than in the
original's advdif recipe --- effectively the $\beta = 0$ ablation.

\textbf{Fix:} raise \code{gamma} to $10^6$ (with D2's rescaling, or rescale
consistently).

\subsection*{D4 --- Augmentation-loss broadcasting bug}

\code{phys\_func(aug\_x\_mean)} outputs the concatenated $(\mu, \log\sigma)$ of shape
$(B, 2)$, but it is compared against \code{aug\_param} of shape $(B, 1)$
(\code{train\_phys\_vae.py:161}). Broadcasting silently regresses \emph{both} $\mu$
and $\log\sigma$ toward the coefficient value. Consequences: $\log\sigma$ is pushed
toward $a \in [0.015, 0.21]$, giving a posterior standard deviation around
$e^{0.1} \approx 1.1$ (astronomically wide for $a \in [0.01, 0.1]$), and every
downstream use of sampled latents (training decode, evaluation) carries huge noise.

\textbf{Fix:} split the output and regress $\mu$ only --- ideally through the same
softplus-style head as the original.

\subsection*{D5 --- Discrepancy loss compares the wrong tensors}

\code{full\_forward} returns x-means reshaped to $(mB, T, L)$ with $m$-major ordering
(\code{phys\_vae.py:390}), but \code{discrepancy\_loss} slices
\code{x\_stats[0][:-1]} / \code{[1:]} as if the first axis were $(m, B, \dots)$
(\code{train\_phys\_vae.py:187--196}). With $m = 2$ the masked entry is
$\|x[\text{batch }0,\text{combo }0] - x[\text{batch }1,\text{combo }0]\|^2$ ---
the distance between \emph{two different data points'} physics-only reconstructions,
not between the same point with and without the auxiliary function. The least-action
regularizer therefore does not do its job (preventing the auxiliary functions from
overriding physics) and actively pulls different samples' outputs together.

\textbf{Fix:} reshape to $(m, B, T, L)$ and take per-sample pairwise combination
differences.

\subsection*{D6 --- Evaluation uses a sampled latent, not the posterior mean}

\code{eval.py:74} and \code{phys\_vae.py:566--577} evaluate a reparameterized draw
($\mu + \varepsilon e^{\log\sigma}$, clipped), so the MAE includes $e^{\log\sigma}$
sampling noise --- which D4 makes enormous. The original protocol reports the
posterior statistic.

\textbf{Fix:} return and evaluate $\mu$ (after softplus/clamp), not a draw.

\subsection*{D7 --- $200\times$ smaller training budget plus protocol bugs}

500 gradient steps versus the original's 100k; no best-validation checkpointing (the
final model is saved); and the validation loop actually trains on the validation data
(\code{src/train.py:133}), contaminating both the metric and the model. The original
also uses gradient clipping 5.0 and weight decay $10^{-6}$, which matter for stability
when an ODE solver sits inside the training loop.

\subsection*{D8 --- Reconstruction-loss scale}

Sum over time only, then mean over batch \emph{and space}
(\code{train\_phys\_vae.py:86}) --- about $12\times$ smaller than the original's
sum-over-all-dims convention, further tilting the reconstruction/KL/regularizer
balance (interacts with D2 and D3).

\subsection*{Secondary deviations (real, but not accuracy-killers alone)}

\begin{itemize}[nosep]
  \item Linear (not softplus) $\mu$ head with hard clipping to the data range
        (dead gradients outside $[0.01, 0.1]$) versus the original's softplus with a
        wider feasible range $[5\cdot10^{-3},\,2\cdot10^{-1}]$.
  \item Prior range equal to the data range versus the original's deliberately wider,
        offset prior.
  \item 12-point grid with the double-\code{jnp.gradient} wide stencil versus a
        20-point grid with the 3-point stencil (weaker curvature signal).
  \item Frozen-boundary versus ghost-cell boundary handling.
  \item Identity decoder map versus a learned linear map.
  \item \code{t1 = integral\_len} (a step count passed as a time) makes every decode
        integrate to $t = 50$ while only saving up to $t = 2.45$ --- about $20\times$
        wasted compute per decode (\code{phys\_vae.py:345--350}).
\end{itemize}

\subsection*{Integration caveat worth verifying in the working tree}

The committed code is internally inconsistent about tensor shapes.
\code{loss\_fn} and the encoder's \code{unmixer\_map} require 3-D batches
$(B, T, L)$ (\code{train\_phys\_vae.py:231}, \code{phys\_vae.py:139}), but the config
never sets \code{retain\_dim: true} for the \code{pivae} dataset (the default
\code{False} flattens to $(B, T\cdot L)$), and the evaluation path
\code{PhysVAE.\_\_call\_\_} slices \code{x[:, :num\_grid\_points]}, which is only
correct for 2-D time-major flattened data (\code{phys\_vae.py:573}). As committed,
training crashes without \code{retain\_dim: true} and evaluation is then
shape-inconsistent. Since accuracy numbers (rather than crashes) are being observed,
the runs presumably use a locally patched variant --- it is worth confirming which
convention the working tree actually uses, because a wrong \code{init\_y} slice in
evaluation would also corrupt the coefficient metric.

\section{Suggested fix order}
\label{sec:fixes}

\begin{enumerate}
  \item Set \code{timestep} to the true data $dt$ ($1/49$, or regenerate the data with
        $dt = 0.05$); make $dt$ flow from the generator into the config, as in the
        original's \code{args.json}.
  \item Set \code{x\_log\_sigma} so that $e^{\,\code{x\_log\_sigma}} \approx 10^{-6}$
        (i.e.\ $\approx -13.8$), restoring the paper's KL balance.
  \item Raise the augmentation weight to $\sim10^6$ and fix its $\mu/\log\sigma$ split
        bug (D4).
  \item Fix the discrepancy loss to compare combinations within a sample (D5).
  \item Evaluate the posterior mean; report relative error for paper-comparable
        numbers (D6).
  \item Extend training ($\ge 5$k epochs at this batch size), stop updating on the
        validation loader, select the best-validation-reconstruction checkpoint, add
        gradient clipping (D7).
  \item Optionally: a softplus $\mu$ head with a feasible range wider than the data
        range, and a wider prior on $a$.
\end{enumerate}

Fixes 1--3 should account for the bulk of the gap: D1 sets the scale of $a$ (the
$2.45\times$ bias), D2 removes the collapse pressure toward the prior mean, and D3
restores the term the paper identifies as decisive for coefficient inference. The
ranking is from static analysis of both codebases; D1 and D2 are verifiable in one
short run by checking whether the predicted $a$ clusters near $0.41\,a_\mathrm{true}$
(indicative of D1) or near the prior mean $0.055$ (indicative of D2).

\end{document}
